\chapter{Risque de taux}\label{ch:rc:rates-risk}

Un desk de swaps en dollars détient deux cents swaps, pour 12~milliards de dollars de notionnel, et son
rapport de risque affiche un DV01 parallèle nul~: une hausse d'un point de base de tous les
taux le laisserait exactement là où il est. Rejouons sur lui les mouvements de la courbe du Trésor
américain du 21 octobre 2022, quand le rendement à deux ans a baissé de 13 points de base et
celui à trente ans a monté de 9, et il perd trois millions de dollars. Le rapport avait raison
sur le niveau et muet sur la forme. Le risque de taux est un vecteur, un nombre par point de la
courbe, et le gérer revient à choisir les coordonnées dans lesquelles lire ce vecteur, à savoir
comment la courbe bouge réellement, et à couvrir un risque réparti sur de nombreux points avec
une poignée d'instruments. Ce chapitre traite ces trois aspects à l'aide du constructeur de courbe
du chapitre~1, avec dix ans de données quotidiennes du Trésor comme modèle de la façon dont les
courbes bougent.

\section{Sensibilités par pilier et échéanciers de risque}

\begin{definition}[Sensibilité par pilier, échéancier de risque]\label{def:rc:rates-risk:ladder}
Une \emph{sensibilité par pilier}\index{sensibilité par pilier} $\Delta_k$ d'une
position est la variation de sa valeur lorsque la $k$-ième cotation d'entrée de la courbe
monte d'un point de base et que la courbe est recalibrée, les autres étant maintenues. Le
\emph{échéancier de risque}\index{échéancier de risque} (risk ladder) est le vecteur $(\Delta_1,\dots,\Delta_n)$
sur l'ensemble des entrées~; sa somme est la sensibilité à un déplacement parallèle des cotations,
et son produit scalaire avec un vecteur de variations de cotations $\delta q$ (en points de base)
est le P\&L au premier ordre, $\pnl\approx\sum_k\Delta_k\,\delta q_k$.
\end{definition}

Un swap au pair qui arrive à échéance sur un \omterm{def:rc:curve-construction:pillar}{pilier} ne charge que ce \omterm{def:rc:curve-construction:pillar}{pilier} sous une
interpolation locale (\cref{prop:rc:curve-construction:triangular})~: sur la courbe à huit \omterm{def:rc:curve-construction:pillar}{piliers} du chapitre
(taux forward plats, swaps SOFR au pair à 1, 2, 3, 5, 7, 10, 20 et
30 ans), un payeur à dix ans de 100~millions de dollars a un échéancier de 83\,778~dollars sur la
cotation à dix ans et zéro ailleurs. L'échéancier d'un livre est la somme de ceux de ses transactions.

\begin{example}[L'échéancier du desk]\label{ex:rc:rates-risk:ladder}
Le livre du chapitre compte 200 swaps d'échéance, de sens et de taille aléatoires (avec germe fixe),
plus un payeur à deux ans de 1.471~milliard de dollars et un payeur à trente ans de
281~millions de dollars ajoutés pour que le DV01 parallèle soit nul. Son échéancier
(\cref{fig:rc:rates-risk:ladder}) est tout sauf nul~: $+212\,010$ à deux ans,
$+249\,817$ à sept, $-279\,791$ à dix, $-1\,135\,816$ à vingt et
$+958\,308$ à trente, par point de base.
\end{example}

\begin{omfigure}
\begin{tikzpicture}
\begin{axis}[width=12cm,height=5.4cm,ybar,bar width=9pt,
  ylabel={milliers de dollars par pb},
  tick label style={font=\small},label style={font=\small},
  xmin=-0.6,xmax=7.6,ymin=-1250,ymax=1100,grid=major,grid style={gray!25},
  xtick=data,xticklabels={1A,2A,3A,5A,7A,10A,20A,30A},
  legend columns=2,legend style={font=\footnotesize,at={(0.5,-0.2)},anchor=north,draw=none,fill=none,/tikz/every even column/.append style={column sep=8pt}},
  table/col sep=comma]
\addplot[fill=omDef!55,draw=omDef] table[x=i,y=book]{figdata/rates-credit-risk/03-rates-risk/ladder.csv};
\addplot[fill=omThm!55,draw=omThm] table[x=i,y=hedged]{figdata/rates-credit-risk/03-rates-risk/ladder.csv};
\legend{livre,livre avec couverture à variance minimale}
\end{axis}
\end{tikzpicture}

\omcaption{L'\omterm{def:rc:rates-risk:ladder}{échéancier de risque} du livre de 200 swaps du chapitre, dont les compartiments
somment à zéro, et l'échéancier après la couverture à trois swaps de
l'\cref{ex:rc:rates-risk:hedge} (swaps à deux, vingt et trente ans). La couverture
n'annule pas ces trois compartiments~; elle laisse la combinaison de risques que les mouvements
historiques rendent la moins coûteuse à porter. Données~: le livre illustratif du chapitre et
le tutoriel.}\label{fig:rc:rates-risk:ladder}
\end{omfigure}

\section{Du risque par pair au risque zéro~: la jacobienne}

L'échéancier ci-dessus est exprimé dans les coordonnées des cotations d'entrée. Les modèles, les
moteurs de scénarios et les régulateurs le veulent souvent dans d'autres coordonnées~: taux zéro
aux \omterm{def:rc:curve-construction:pillar}{piliers}, taux forward, ou un ensemble standard d'échéances clés.

\begin{definition}[Jacobienne de la courbe]\label{def:rc:rates-risk:jacobian}
La \emph{jacobienne de la courbe}\index{jacobienne de la courbe} est la matrice $J_{kj} =
\partial m_k/\partial z_j$ des dérivées des cotations modèles des instruments d'entrée par rapport
aux paramètres aux \omterm{def:rc:curve-construction:pillar}{piliers} de la courbe, sur la courbe calibrée~; son inverse donne la façon dont
les \omterm{def:rc:curve-construction:pillar}{piliers} bougent lorsque les cotations bougent, $\delta z
= J^{-1}\delta q$.
\end{definition}

\begin{proposition}[Changement de coordonnées]\label{prop:rc:rates-risk:jacobian}
Soient $g_q = (\partial V/\partial q_k)_k$ et $g_z = (\partial V/\partial
z_j)_j$ les échéanciers d'une position en cotations et en taux zéro aux \omterm{def:rc:curve-construction:pillar}{piliers}. Alors
\[
g_z = J^\top g_q,\qquad g_q = (J^\top)^{-1}g_z .
\]
\end{proposition}
\begin{proof}
La position ne dépend des cotations qu'à travers les \omterm{def:rc:curve-construction:pillar}{piliers}, $V = V(z(q))$, donc
$g_q = (\partial z/\partial q)^\top g_z = (J^{-1})^\top g_z$.
\end{proof}

\begin{omfigure}
\begin{tikzpicture}[font=\small,
  box/.style={draw,thick,rounded corners=2pt,align=center,minimum height=1cm,minimum width=2.6cm,inner sep=4pt},
  arr/.style={-{Stealth[length=2.5mm]},thick}]
\node[box,draw=omDef,fill=omDef!8] (q) at (0,0) {cotations d'entrée\\$q_1,\dots,q_n$};
\node[box,draw=omProp,fill=omProp!8] (z) at (5,0) {taux zéro aux piliers\\$z_1,\dots,z_n$};
\node[box,draw=omThm,fill=omThm!8] (v) at (10,0) {valeur de la\\position $V$};
\draw[arr] (q) -- node[above=4pt] {calibration} node[below=4pt] {$\delta z = J^{-1}\delta q$} (z);
\draw[arr] (z) -- node[above=4pt] {valorisation} node[below=4pt] {$g_z=\partial V/\partial z$} (v);
\draw[arr,omExo,dashed] (v.south) to[bend left=22] node[below=3pt] {risque en cotations~: $g_q = (J^\top)^{-1}g_z$} (q.south);
\end{tikzpicture}

\omcaption{Deux systèmes de coordonnées pour un même risque. La calibration fait correspondre
cotations et \omterm{def:rc:curve-construction:pillar}{piliers}, la valorisation fait correspondre \omterm{def:rc:curve-construction:pillar}{piliers} et valeur~; la jacobienne de la
calibration fait passer un échéancier d'un système à l'autre
(\cref{prop:rc:rates-risk:jacobian}).}\label{fig:rc:rates-risk:jacobian}
\end{omfigure}

\begin{example}[Deux échéanciers d'un même swap]\label{ex:rc:rates-risk:twoladders}
L'échéancier par pair du payeur à dix ans vaut $83\,778$ sur la seule cotation à dix ans~; son
échéancier en taux zéro vaut $367$, $713$, $1\,531$, $3\,205$, $5\,120$ et $74\,016$ sur les
\omterm{def:rc:curve-construction:pillar}{piliers} d'un à dix ans. Relever un taux zéro précoce diminue la valeur des coupons fixes que
le payeur doit, et dans l'espace des cotations cet effet est absorbé par la recalibration. Pour
l'ensemble du livre, l'échéancier par pair somme à zéro et l'échéancier en taux zéro à
11\,408~dollars~: un déplacement parallèle des cotations et un déplacement parallèle des taux zéro
sont deux scénarios différents.
\end{example}

\begin{definition}[Duration partielle]\label{def:rc:rates-risk:krd}
Une \emph{duration partielle}\index{duration partielle} (key-rate duration) (Ho, 1992) est la
sensibilité d'une position à un déplacement en «~tente~» de la courbe des taux zéro
centré sur une maturité clé~: un point de base à la clé, décroissant linéairement
jusqu'à zéro aux deux clés voisines. Les tentes somment à un à chaque échéance, de
sorte que les sensibilités par duration partielle s'additionnent pour donner la
sensibilité à un déplacement parallèle des taux zéro~; elles ne dépendent pas de la
manière dont la courbe a été construite, ce qui en fait la monnaie habituelle des rapports de risque entre desks.
\end{definition}

\section{Composantes principales de la courbe}

Comment la courbe se déplace-t-elle~? La matrice de covariance des variations
quotidiennes y répond, et ses vecteurs propres, les composantes principales du livre~4,
chapitre~22, lui donnent une forme.

\begin{definition}[Facteurs de niveau, de pente et de courbure]\label{def:rc:rates-risk:factors}
Pour une courbe échantillonnée en $n$ maturités, soient $v_1, v_2, v_3$ les trois
premiers vecteurs propres de la covariance de ses variations quotidiennes (avec $\lambda_1\ge
\lambda_2\ge\lambda_3$). Le \emph{facteur de niveau}\index{facteur de niveau} (level factor) $v_1$ a des
charges de même signe, le \emph{facteur de pente}\index{facteur de pente} (slope factor) $v_2$
change de signe une fois, des maturités courtes aux maturités longues, et le \emph{facteur de
courbure}\index{facteur de courbure} (curvature factor) $v_3$ change de signe deux fois. Le mouvement du jour $t$ a
pour score $v_i^\top\delta q_t$ sur le facteur $i$, et les facteurs sont non corrélés
par construction.
\end{definition}

\begin{example}[Dix ans de mouvements du Trésor américain]\label{ex:rc:rates-risk:pca}
Sur les 2\,681 variations quotidiennes de la courbe par pair du Trésor américain de janvier 2016
à septembre 2026 (huit maturités, de un à trente ans), les trois
premières composantes expliquent 85.1~\%, 11.0~\% et 2.1~\% de la variance, soit 98.2~\% au total,
avec des écarts-types quotidiens de 13.6, 4.9 et 2.1 points de base de score
(\cref{fig:rc:rates-risk:pca}). Le niveau pèse moins sur la maturité d'un an que sur le
ventre~; le \omterm{def:rc:rates-risk:factors}{facteur de courbure} est dominé par le rendement à un an, la maturité la plus
liée aux prochaines réunions de politique monétaire. Litterman et Scheinkman ont retrouvé les
mêmes trois facteurs en 1991.
\end{example}

\begin{omfigure}
\begin{tikzpicture}
\begin{axis}[width=11.5cm,height=5.6cm,
  xlabel={maturité (années, échelle log)},ylabel={charge},xmode=log,
  xtick={1,2,3,5,7,10,20,30},xticklabels={1,2,3,5,7,10,20,30},
  tick label style={font=\small},label style={font=\small},
  xmin=0.9,xmax=33,ymin=-0.6,ymax=0.9,grid=major,grid style={gray!25},
  legend columns=3,legend style={font=\footnotesize,at={(0.5,-0.25)},anchor=north,draw=none,fill=none,/tikz/every even column/.append style={column sep=8pt}},
  table/col sep=comma]
\addplot[omDef,very thick,mark=*] table[x=m,y=pc1]{figdata/rates-credit-risk/03-rates-risk/pca.csv};
\addplot[omThm,thick,mark=square*] table[x=m,y=pc2]{figdata/rates-credit-risk/03-rates-risk/pca.csv};
\addplot[omProp,thick,dashed,mark=triangle*] table[x=m,y=pc3]{figdata/rates-credit-risk/03-rates-risk/pca.csv};
\legend{niveau (85.1~\%),pente (11.0~\%),courbure (2.1~\%)}
\end{axis}
\end{tikzpicture}

\omcaption{Charges des trois premières composantes principales des variations quotidiennes
des taux par pair du Trésor américain, de janvier 2016 à septembre 2026, avec leur part
de variance. Source~: US Treasury, courbe quotidienne des taux par pair (domaine public)~;
le tutoriel du chapitre.}\label{fig:rc:rates-risk:pca}
\end{omfigure}

\begin{remark}[La perte du hook, facteur par facteur]\label{rem:rc:rates-risk:decomp}
En décomposant le mouvement du 21 octobre 2022 et l'échéancier du livre sur les
huit composantes, le \omterm{def:rc:rates-risk:factors}{facteur de pente} (score $+21.2$) explique $-\text{USD}~3.09$
millions sur la perte du premier ordre de $-3.00$ millions, et les composantes d'ordre supérieur
ajoutent $+0.72$, $+0.69$ et $-1.48$ millions, qui se compensent presque. Le livre était
neutre au niveau et short d'un steepener.
\end{remark}

\section{Couvrir un livre de swaps}

L'échéancier d'un livre compte autant de paniers que la courbe compte de \omterm{def:rc:curve-construction:pillar}{piliers}~; un desk le
couvre avec quelques swaps liquides. Lesquels, et en quelle quantité chacun, dépend de la
manière dont les paniers évoluent ensemble.

\begin{method}[Couverture à variance minimale]\label{met:rc:rates-risk:mvhedge}
Soit $g$ l'échéancier du livre, $H$ la matrice dont les colonnes sont les échéanciers
d'une unité de chaque instrument de couverture, et $\Sigma$ la matrice de covariance des
variations quotidiennes des cotations. Le P\&L quotidien du livre plus la couverture vaut $(g+Hh)^\top\delta q$, de variance
$(g+Hh)^\top\Sigma(g+Hh)$, minimisée par
\[
h^\star = -\bigl(H^\top\Sigma H\bigr)^{-1}H^\top\Sigma g .
\]
Indiquer la part résiduelle de variance $(g+Hh^\star)^\top\Sigma(g+Hh^\star) /
g^\top\Sigma g$, et choisir les instruments qui la rendent la plus petite.
\end{method}

\omcode{code/firm/ratesrisk/firm_ratesrisk.py}{94}{101}{La couverture à variance minimale et la variance qu'elle laisse.}\label{lst:rc:rates-risk:hedge}

\begin{example}[Trois swaps contre deux cents]\label{ex:rc:rates-risk:hedge}
En prenant les variations des taux par pair du Trésor comme substituts des variations
des taux de swap, le P\&L quotidien du livre a un écart-type de USD~1\,288\,843. Le trio
classique de swaps à deux, dix et trente ans laisse 58.3~\% de la variance~; couvrir exactement les trois
composantes principales en laisse 68.9~\%, car l'essentiel du risque du livre
réside dans le spread vingt-trente ans, que les trois premiers facteurs ne décrivent
pas. Les trois meilleurs swaps sont ceux à deux, vingt et trente ans~:
recevoir le fixe sur USD~1.196 milliard à deux ans, payer sur USD~801 millions à
vingt ans et recevoir sur USD~515 millions à trente ans. Ils ne laissent que 4.2~\%
de la variance, soit un écart-type quotidien de USD~264\,554
(\cref{fig:rc:rates-risk:pnl}).
\end{example}

\begin{omfigure}
\begin{tikzpicture}
\begin{axis}[width=12cm,height=5.4cm,
  xlabel={date},ylabel={P\&L quotidien (millions de USD)},
  xtick={2025,2025.5,2026,2026.5},xticklabels={janv. 2025,juil. 2025,janv. 2026,juil. 2026},
  tick label style={font=\small},label style={font=\small},
  xmin=2024.66,xmax=2026.75,ymin=-6,ymax=4.5,grid=major,grid style={gray!25},
  legend columns=2,legend style={font=\footnotesize,at={(0.5,-0.25)},anchor=north,draw=none,fill=none,/tikz/every even column/.append style={column sep=8pt}},
  table/col sep=comma]
\addplot[omDef!70,thin] table[x=t,y=book]{figdata/rates-credit-risk/03-rates-risk/pnl.csv};
\addplot[omThm,thick] table[x=t,y=hedged]{figdata/rates-credit-risk/03-rates-risk/pnl.csv};
\legend{livre (premier ordre),livre avec la couverture à trois swaps}
\end{axis}
\end{tikzpicture}

\omcaption{P\&L quotidien de premier ordre du livre, $g^\top\delta q_t$, sur
les mouvements des taux au pair du Trésor américain de septembre 2024 à
septembre 2026, sans couverture et avec la couverture de variance minimale à
deux, vingt et trente ans (estimée sur 2016--2026). Source~: taux au pair du
Trésor américain~; tutoriel du chapitre.}\label{fig:rc:rates-risk:pnl}
\end{omfigure}

\begin{remark}[Ce que suppose la couverture]\label{rem:rc:rates-risk:assumes}
La couverture ne vaut que par sa covariance~: dix ans de mouvements du Trésor
incluent les années de taux nuls et les hausses de 2022, et les taux de swap
évoluent avec les taux du Trésor plus les spreads de swap, ce que ce modèle ignore. Une couverture estimée
sur un régime et utilisée dans un autre laisse plus de 4~\%~; les tests de résistance
(\cref{ch:rc:stress-testing-and-scenarios}) existent précisément pour cela.
\end{remark}

\section{Gamma et gamma croisé}

L'échéancier est une dérivée première~; pour les grands mouvements, les dérivées
secondes comptent.

\begin{definition}[Gamma croisé]\label{def:rc:rates-risk:crossgamma}
Le \emph{gamma croisé}\index{gamma croisé} d'une position entre les facteurs $k$ et
$l$ est $\Gamma_{kl} = \partial^2V/\partial q_k\partial q_l$ pour $k\ne l$, par
point de base au carré~; avec les termes diagonaux, il forme la matrice de gamma, et le
P\&L du second ordre vaut $\sum_k\Delta_k\delta q_k +
\tfrac12\sum_{k,l}\Gamma_{kl}\delta q_k\delta q_l$.
\end{definition}

\begin{example}[Le terme du second ordre du crochet]\label{ex:rc:rates-risk:gamma}
Sur le mouvement du 21 octobre 2022, l'échéancier du livre donne une perte de
premier ordre exactement égale à 3\,000\,000~USD (par construction)~; la réévaluation complète fournit
$-3\,077\,970$. La matrice de gamma, calculée par différences centrées, donne
$\tfrac12\,\delta q^\top\Gamma\,\delta q = -78\,189$, ce qui explique l'écart
à 220~USD près. Le gamma parallèle du livre vaut $-377$ dollars par
point de base au carré~: il est vendeur de convexité, surtout à travers ses
swaps payeurs à trente ans.
\end{example}

Pour un livre de swaps, le \omterm{def:rc:rates-risk:crossgamma}{gamma croisé} est faible et provient surtout de
l'actualisation~; pour les options, et pour les produits dont le paiement dépend
de deux taux (chapitre~6), il est de premier ordre.

\section{Tutoriel~: échelle, facteurs et couverture}\label{tut:rc:rates-risk:tutorial}

\begin{tutorial}
\textbf{Objectif.} Mesurer le risque du livre en coordonnées de cotation, de taux
zéro et de \omterm{def:rc:rates-risk:krd}{duration partielle}, trouver comment évolue la courbe du Trésor, et
couvrir le livre avec trois swaps. \textbf{État final~:}
\cref{fig:rc:rates-risk:ladder,fig:rc:rates-risk:pca,fig:rc:rates-risk:pnl} et les
nombres de l'\cref{ex:rc:rates-risk:hedge}.
\begin{tutsteps}
\item \textbf{Le livre}~: \texttt{balanced\_book()} renvoie 202 swaps dont le DV01
parallèle est nul.
\item \textbf{Échelles}~: \texttt{zero\_and\_key\_ladders(book)}~; vérifier que
\texttt{par\_to\_zero} de l'échelle des taux au pair égale l'échelle des taux zéro.
\omcode{code/firm/ratesrisk/firm_ratesrisk.py}{19}{27}{L'échelle des taux au pair~: choquer chaque cotation, recalibrer, réévaluer.}\label{lst:rc:rates-risk:ladder}
\item \textbf{Facteurs}~: \texttt{components()} sur le fichier du Trésor~; les parts
doivent valoir 85.1, 11.0 et 2.1~\%.
\item \textbf{Couverture}~: \texttt{hedge(book)} et \texttt{hedge(book, (1, 5, 7))}
pour le trio deux, dix et trente ans~; \texttt{fig\_rc\_ratesrisk.py} écrit
les graphiques.
\end{tutsteps}
\textbf{Pour aller plus loin.} Estimer la covariance sur 2016--2020 seulement et
mesurer la variance couverte sur 2021--2026~; remplacer les taux forward plats par
des taux monotones convexes et comparer les échelles.
\end{tutorial}

\section{Construction~: la boîte à outils de transformation du risque}\label{bld:rc:rates-risk:ratesrisk}

\begin{build}
\bfield{Rôle} La couche de risque de taux de la mini-firme~: l'échelle de chaque
position de taux, dans n'importe quelles coordonnées, et la couverture qui la
neutralise~; le moteur de risque du \cref{ch:rc:build-a-risk-engine} l'appelle.
\bfield{Interface} \texttt{par\_ladder(spot, instruments, kind, pv)},
\texttt{zero\_ladder(curve, pv)}, \texttt{par\_to\_zero}, \texttt{zero\_to\_par},
\texttt{TentBumped}, \texttt{key\_rate\_ladder}, \texttt{pca(changes)},
\texttt{min\_variance\_hedge(g, H, cov)}, \texttt{cross\_gamma}.
\bfield{Règles} Sensibilités par point de base à partir de petits chocs mis à
l'échelle (non-linéarité du chapitre~1)~; signes de l'ACP fixés (niveau positif,
pente croissante avec la maturité)~; aucun état global.
\bfield{Tests d'acceptation} \texttt{code/firm/ratesrisk/tests/}~: le jacobien
envoie les échelles au pair sur les échelles zéro pour deux interpolations~; les sensibilités par \omterm{def:rc:rates-risk:krd}{duration partielle}
somment à la sensibilité zéro parallèle~; l'ACP retrouve une structure à deux facteurs
plantée~; couvrir un livre par lui-même ne laisse rien~; la matrice de gamma est
symétrique et explique le P\&L du second ordre d'un swap à trente ans.
\bfield{Pour aller plus loin} Compartiments de taux forward~; un jacobien
analytique à partir du solveur du chapitre~1~; une couverture avec coûts de transaction (une pénalité sur $|h|$).
\end{build}

\begin{omsources}
\item R.\ Litterman et J.\ Scheinkman, ``Common factors affecting bond
returns'', \emph{Journal of Fixed Income} 1(1), 1991.
\item T.\ S.\ Y.\ Ho, ``Key rate durations: measures of interest rate risks'',
\emph{Journal of Fixed Income} 2(2), 1992.
\item US Department of the Treasury, \emph{Daily Treasury Par Yield Curve Rates}.
\item L.\ Andersen et V.\ Piterbarg, \emph{Interest Rate Modeling}, volume 3,
Atlantic Financial Press, 2010, chapitre 26 (couverture en vega et en delta des livres de taux).
\end{omsources}

\section{Exercices}

\begin{exercise}[$\star$]\label{exo:rc:rates-risk:1}
Un livre a pour échelle $+50$, $-120$, $+80$ milliers de dollars par point de base
sur les cotations à deux, dix et trente ans. Calculer son DV01 parallèle et son
P\&L si les trois cotations bougent de $-5$, $+2$ et $+6$ points de base.
\end{exercise}

\begin{exercise}[$\star$]\label{exo:rc:rates-risk:2}
Pourquoi l'échelle des taux au pair de ce livre somme-t-elle à zéro alors que son échelle des taux zéro somme à
11\,408~USD~?
\end{exercise}

\begin{exercise}[$\star$]\label{exo:rc:rates-risk:3}
En reprenant l'\cref{ex:rc:rates-risk:pca}, quelle part de la variance quotidienne subsiste après
les trois premières composantes, et quelle est l'amplitude d'un mouvement
quotidien d'un écart-type du \omterm{def:rc:rates-risk:factors}{facteur de niveau}, en points de base de score~?
\end{exercise}

\begin{exercise}[$\star\star$]\label{exo:rc:rates-risk:4}
Calculer, à partir de l'\cref{ex:rc:rates-risk:twoladders}, la somme de l'échelle des taux zéro du payeur à dix ans
et la comparer à son échelle des taux au pair. Laquelle est la plus grande, et pourquoi~?
\end{exercise}

\begin{exercise}[$\star\star$]\label{exo:rc:rates-risk:5}
Pourquoi une échelle des \omterm{def:rc:rates-risk:krd}{durations partielles} est-elle indépendante de l'interpolation de la courbe,
alors qu'une échelle au pair ne l'est pas. Laquelle transmettre à un
régulateur, et laquelle à un trader qui couvrira avec des swaps au pair~?
\end{exercise}

\begin{exercise}[$\star\star$]\label{exo:rc:rates-risk:6}
Lire la \cref{fig:rc:rates-risk:ladder}. Après la couverture à variance minimale, les
\omterm{def:rc:rates-risk:ladder}{sensibilités par pilier} à sept et dix ans restent inchangées à $+250$ et $-280$ milliers. Pourquoi
la couverture les a-t-elle laissées, et qu'ont-elles en commun~?
\end{exercise}

\begin{exercise}[$\star\star\star$]\label{exo:rc:rates-risk:7}
\emph{Programmation.} Trouver la meilleure couverture à un seul swap et la meilleure
couverture à quatre swaps du livre, ainsi que la variance que chacune laisse.
\end{exercise}

\begin{exercise}[$\star\star\star$]\label{exo:rc:rates-risk:8}
\emph{Trouver la faille.} «~Nous neutralisons les \omterm{def:rc:rates-risk:factors}{facteurs de niveau}, de pente et de courbure
chaque soir, donc nous ne portons aucun risque de courbe.~» Corriger-la à l'aide des chiffres du livre.
\end{exercise}

\section{Problème~: La journée tordue}

\begin{problem}[{Problème du week-end --- un livre neutre en DV01 le 21 octobre 2022}]\label{pb:rc:rates-risk:1}
Le livre du chapitre (202 swaps, DV01 parallèle nul) est rejoué sur les variations des
taux au pair du 21 octobre 2022~: $-8$, $-13$, $-14$, $-11$, $-8$, $-3$, $+7$, $+9$
points de base, de un à trente ans.

\medskip\noindent\textbf{Partie I --- La perte.}
\begin{enumerate}
\item Donner le P\&L du premier ordre à partir de l'échelle.
\item Donner le P\&L par réévaluation complète, et l'écart.
\item Quelles sont les deux \omterm{def:rc:rates-risk:ladder}{sensibilités par pilier} qui contribuent le plus, avec leurs signes~?
\item Le mouvement était-il une pentification ou un aplatissement, et dans quel sens le livre
était-il positionné~?
\item Pourquoi le DV01 parallèle n'a-t-il rien signalé~?
\end{enumerate}

\medskip\noindent\textbf{Partie II --- Facteurs.}
\begin{enumerate}[resume]
\item Donner le score du mouvement sur le \omterm{def:rc:rates-risk:factors}{facteur de pente}.
\item Donner la contribution du \omterm{def:rc:rates-risk:factors}{facteur de pente} à la perte.
\item Combien d'écarts types représentait ce mouvement de pente~?
\item Les composantes d'ordre supérieur contribuent $+0.72$, $+0.69$ et $-1.48$ million. Que
représentent-elles~?
\item Un modèle factoriel à trois facteurs aurait-il prévu la perte~?
\end{enumerate}

\medskip\noindent\textbf{Partie III --- Couverture.}
\begin{enumerate}[resume]
\item Donner l'écart type quotidien du P\&L non couvert.
\item Donner la part résiduelle obtenue avec des swaps à deux, dix et trente ans.
\item Donner les trois meilleurs swaps, leurs notionnels et leurs sens.
\item Donner l'écart type quotidien résiduel avec eux.
\item Donner le P\&L du premier ordre du livre couvert le 21 octobre 2022.
\end{enumerate}

\medskip\noindent\textbf{Partie IV --- Jugement.}
\begin{enumerate}[resume]
\item Pourquoi ne pas simplement couvrir chaque \omterm{def:rc:rates-risk:ladder}{sensibilité par pilier} à zéro~?
\item À quelle fréquence la covariance doit-elle être réestimée, et sur quelle fenêtre~?
\item Que faudrait-il ajouter au rapport de risque quotidien~?
\item Énoncer le \emph{résultat nommé}~: la perte du jour tordu, et les trois
swaps qui suppriment l'essentiel du risque, ainsi que la variance qu'ils laissent.
\item En une phrase~: pourquoi le DV01 n'est-il pas une mesure de risque~?
\end{enumerate}
\end{problem}

\section{Questions d'entretien}

\begin{interviewq}[$\star$ \iqroles{trader, risque}]\label{iq:rc:rates-risk:1}
Votre livre a un DV01 nul. Êtes-vous couvert~?
\end{interviewq}

\begin{interviewq}[$\star\star$ \iqroles{chercheur, développeur}]\label{iq:rc:rates-risk:2}
Comment convertissez-vous un risque exprimé en cotations de swaps au pair en risque sur taux zéro, et inversement~?
\end{interviewq}

\begin{interviewq}[$\star\star$ \iqroles{chercheur, trader}]\label{iq:rc:rates-risk:3}
Quelles sont les trois premières composantes principales des variations de la courbe des taux, et
quelle part de la variance expliquent-elles~?
\end{interviewq}

\begin{interviewq}[$\star\star$ \iqroles{trader}]\label{iq:rc:rates-risk:4}
Vous devez couvrir un livre avec seulement trois swaps. Comment les choisissez-vous~?
\end{interviewq}

\begin{interviewq}[$\star\star\star$ \iqroles{risque, chercheur}]\label{iq:rc:rates-risk:5}
Votre explication du P\&L au premier ordre se trompe de 3~\% un jour de grand mouvement. Qu'est-ce
qui manque, et comment le calculez-vous~?
\end{interviewq}

\begin{interviewq}[$\star\star\star$ \iqroles{chercheur}]\label{iq:rc:rates-risk:6}
Dans quels cas ne feriez-vous pas confiance à une couverture fondée sur une ACP~?
\end{interviewq}
